\documentclass{article}
\usepackage[UTF8]{ctex}
\usepackage{amsmath,amssymb,amsfonts}
\usepackage{enumitem}
\usepackage{geometry}
\geometry{a4paper, margin=1in}
% === 错题本排版优化（不参与题目标识） ===
\setlength{\parindent}{0pt}
\setlength{\parskip}{0.35em plus 0.10em minus 0.08em}
\setlength{\emergencystretch}{3em}
\setlist{itemsep=0.25em, topsep=0.40em, parsep=0pt}
\allowdisplaybreaks[2]
% === 错题本排版优化结束 ===
\begin{document}
\section*{上海财经大学2026年数学分析真题}
\section*{52．上海财经大学2026年数学分析试卷}

\subsection{题目1}
1．判断下述命题是否正确并给出理由。

\begin{itemize}
  \item[（1）] 设数列 \(\left\{b_{n}-a_{n}\right\}\) 无界，数列 \(\left\{a_{n}\right\}\) 收敛，则数列 \(\left\{b_{n}\right\}\) 无界。
  \item[（2）] 设数列 \(\left\{a_{n}\right\}\) 满足 \(\lim _{n \rightarrow \infty} a_{n}=a\) ，且 \(a_{n}\) 是集合 \(A\) 的上界，则 \(a\) 是集合 \(A\) 的上界．

\subsection{题目2}
\end{itemize}

\subsection{题目2}
原题边界“2．”：求极限 \(\lim _{n \rightarrow \infty} \frac{1 \times 3 \times 5 \times \cdots \times(2 n-1)}{2 \times 4 \times 6 \times \cdots \times(2 n)}\) ．

\subsection{题目3}
3．计算 \(\int_{0}^{+\infty} \frac{1}{\left(1+x^{2}\right)\left(1+x^{2026}\right)} \mathrm{d} x\) ．

\subsection{题目4}
4．计算 \(\iint_{x^{2}+y^{2} \leq 1}|3 x+4 y| \mathrm{d} x \mathrm{~d} y\) ．

\subsection{题目5}
5．计算曲线积分 \(\int_{C} \frac{y \mathrm{~d} x-x \mathrm{~d} y}{x^{2}+y^{2}}\) ，其中 \(C\) 是 \(\frac{x^{2}}{4}+\frac{y^{2}}{9}=1\) 从 \((2,0)\) 到（ \(-2,0\) ）的一段曲线，取逆时针方向．

\subsection{题目6}
6．求平面 \(x+y+z=0\) 与椭球面 \(x^{2}+y^{2}+4 z^{2}=1\) 相交而成的椭圆面积．

\subsection{题目7}
7．设

\[
f(x)= \begin{cases}p, & x=\frac{p}{q}\left(p \in \mathbb{N}^{+}, q \in \mathbb{Z} \backslash\{0\}, p, q \text { 互素 }\right), \\ 1, & x=0, \\ 0, & x \text { 是无理数. }\end{cases}
\]

证明：\(f(x)\) 在任意区间上无界．

\subsection{题目8}
8．设 \(f(x)\) 在 \([a, b]\) 上有定义，证明：\(f(x)\) 在 \([a, b]\) 上一致连续的充要条件是 \(f(x)\) 将 \([a, b]\) 上的任意Cauchy 数列映射为 Cauchy 数列。

\subsection{题目9}
9．已知 \(f(x)\) 在 \([a, b]\) 上二阶可导，且有 \(f(a)=f(b)=0\) ，若存在 \(c \in(a, b)\) 使得 \(f(c)>0\) ．证明：存在 \(\xi \in(a, b)\) ，使得 \(f^{\prime \prime}(\xi)<0\) ．

\subsection{题目10}
10．设 \(f(x)\) 在 \([a,+\infty)\) 上三阶可导，证明：若 \(\lim _{x \rightarrow+\infty} f(x), ~ \lim _{x \rightarrow+\infty} f^{\prime \prime \prime}(x)\) 都存在且有限，则

\[
\lim _{x \rightarrow+\infty} f^{\prime}(x)=\lim _{x \rightarrow+\infty} f^{\prime \prime}(x)=\lim _{x \rightarrow+\infty} f^{\prime \prime \prime}(x)=0 .
\]

\subsection{题目11}
11．设

\[
f(x, y)= \begin{cases}x y \cos \frac{1}{\sqrt{x^{2}+y^{2}}}, & (x, y) \neq(0,0) \\ 0, & (x, y)=(0,0)\end{cases}
\]

证明：\(f(x, y)\) 在 \((0,0)\) 处可微，但偏导函数不连续．

\subsection{题目12}
12．设 \(a_{n} \geq 0\) ，且幂级数 \(\sum_{n=1}^{\infty} a_{n} x^{n}\) 的收敛半径为 \(R\) ，证明： \(\lim _{x \rightarrow R^{-}} \sum_{n=1}^{\infty} a_{n} x^{n}=\sum_{n=1}^{\infty} a_{n} R^{n}\) ．
\end{document}
